\section{The five-year proof: timing, ties, the renewal, and per-class extremes}\label{supp:proof}

The five-year proof is close to a fixed appointment rather than a hazard
spread over years: 97\% of the cancellations that fall in the 4.0--8.5
year window post between registration age 6.5 and 7.0. Since the
cancellation record ends on 2 April 2026, a cohort is fully observable
only once it has reached seven years, which makes 2018 the last complete
cohort. On the 3.15 million registrations whose ninth anniversary is
inside the record --- for which failure is settled and no window is
needed --- the contrast is $+2.71$pp within class and registration year,
against $+2.65$pp on the paper's 4.0--8.5 window; conditional on
failing, the probability of failing late carries a contrast of
$-0.01$pp ($t = -0.2$). Response speed does not predict the outcome:
across the 1.22 million scored registrations with a usable
office-action/response pair, passage varies between 46.1\% and 48.1\%
across response-speed quintiles, from a median of seven days to 183, and
the lead penalty is present in both halves ($-3.48 \pm 0.40$ and
$-2.36 \pm 0.40$pp on passing, against $-2.94$ pooled). The restriction
drops 61\% of the sample, disproportionately applications that never
drew a refusal.

Tied scores: within the class$\times$registration-year cell the quintiles
are cut in, $17.1$\% of scored registrations share a lead value with at
least one other filing. Filings that share a class, a filing date and a
goods/services description have identical theme mixes scored against
identical references. Cut points are fixed by sorting on the score and
then on the serial number; sorting ties the opposite way moves the
top-versus-bottom contrast by $0.001$pp, and five pseudorandom orderings
move it by at most $0.002$pp, against a contrast of $2.287$pp.

At the year-ten renewal, fitted over 872{,}079 five-year survivors with a
resolved status, a straight line through renewal against lead explains
1\% of the variation across forty bins and a V explains 89\%
($\Delta$AIC $= 83$); renewal peaks at 64.2\% near $z = -0.4$ and falls
to 45.5\% at the most lagging bin and 50.6\% at the most leading.

Per class, the largest raw penalties sit in tobacco ($+26$pp, the vaping
era), hand tools ($+10$pp), processed foods ($+10$pp) and medical
apparatus ($+9$pp), with negative lifts in beer and soft drinks,
transport and construction; these are the extremes of a 44-class screen,
sitting away from zero partly by selection. Every class's profile is in
the online appendix.

\section{Waves: the theme tour, internet-bearing filings, and surges}\label{supp:waves}

\subsection{The theme tour}

In 1995--1999 the computer, software, information and data services theme
failed at 64.9\% against 47.5--59.0\% in the other large themes. In
2000--2004 it held about 11\% of the leading fifth against 40\% of
filings; the leading language had moved to video, music and games on
electronic media (31\% of the leading fifth against 12\% of filings), and
inside several themes the lagging fifth failed more than the leading one:
$-3.6$pp in electronic media, $-5.5$ in retail store services, $-6.9$ in
apparatus and control systems. From 2008 the leading fifth
over-represented downloadable and online software 15\% against 4\%; the
theme's within-theme contrast in the 2008--2014 era is $+3.3$pp, and the
computer-services theme's own contrast is positive in every era ($+3.6$
to $+5.9$pp) and never reverses.

\subsection{Internet-bearing filings by class}

\begin{table}[!htbp]\centering\footnotesize
\caption{Internet-bearing filings against the rest of their Nice class
(Section~\ref{sec:internet}). Registration is the share of class-records
filed 1995--2018 that reached registration; failure is the share of
unique registrations 2002--2018 cancelled for non-use at age 4.0--8.5
years. Classes with fewer than 500 internet-bearing registrations are
omitted.}
\label{tab:internetS}
\input{results/internet_breakout.tex}
\end{table}

The internet gap in event time (Figure~\ref{fig:webshape}): dating the
internet's arrival in each class as the first filing year the pattern
reaches 2\% of the class's filings (1994--1995 in the technology and
transport classes, 1999--2000 in apparel and chemicals, 2008 in
machinery), the gap between internet-bearing registrations and the rest
starts near $+7$ points, declines to about zero nine to eleven years
after arrival, and settles between $+1$ and $+4$; across the eight
classes observable in both phases the mean gap is $+4.4$ points in years
zero to four and $+1.6$ from year ten, shrinking in six of the eight.
Calendar time disguises this: transport (arrival 1995) peaks at $+16$
points in its 2000 cohort; machinery (arrival 2008) peaks at $+17$ in its
2012--2015 cohorts. Clothing is the standing exception: its
internet-bearing registrations outlast the class in every cohort after
2004.

\begin{figure}[!htbp]\centering
\includegraphics[width=\textwidth]{results/internet_convergence.png}
\caption{The internet gap in failure at the five-year proof, inside each
class. \emph{Left:} internet-bearing minus other registrations' failure
rate, by registration cohort in years since the internet pattern reached
2\% of the class's filings, averaged across classes (band: one standard
deviation across classes). \emph{Right:} five classes in calendar time.
Cells with fewer than 100 registrations in either group are dropped.}
\label{fig:webshape}
\end{figure}

\subsection{Surges}

Measured against the whole corpus at a 1\% share, five themes ever surge
and one clears a 1{,}000-registration floor --- online software, 1999,
whose 6{,}339 in-window registrations failed at 63.7\%
(Figure~\ref{fig:surgesS}). Within-class episodes differ in sign among
themselves: the education-services surge in advertising and retail after
2002 failed $13.7$ points \emph{less} than its class; the
charitable-fundraising surge in finance after 2011 failed $4.0$ points
more; of the episodes large enough to carry their own lead contrast, one
is positive ($+5.8$pp, $t = 2.2$) and the rest indistinguishable from
zero. On the 61 episodes with at least 300 joining registrations over a
six-year entry window (64{,}105 of 86{,}447 joining registrations):
joiners in the surge year fail $-0.8$ points relative to their wave's
mean and joiners five years in $+0.0$, with no gradient between (SEs
$\sim 0.5$); excess failure is $+2.7$, $+3.9$ and $-0.1$ points across
wave-size terciles, correlating with size at $-0.22$; and waves whose
theme share is still at or above its surge level five years on carry
$+4.3$ points of excess against $+1.0$ for receding waves.

\begin{figure}[!htbp]\centering
\includegraphics[width=\textwidth]{results/fig_surges.png}
\caption{Four surge episodes under the 500-theme model. \emph{Left:} the
theme's share of its class's filings by filing year (for the online
software episode, of all filings corpus-wide), on the 25\% systematic
sample the surge screen uses; the open circle marks the surge year.
\emph{Right:} failure at the five-year proof for registrations whose
dominant theme was the surging one, filed in the three-year surge window
(colored, 95\% intervals), against the rest of the same class in the same
cohorts (grey).}
\label{fig:surgesS}
\end{figure}

\section{Success: unfunded listings, value tiers, and curated vocabularies}\label{supp:success}

Of the 1{,}669 debut owners with an IPO marker, 759 have no post-debut
Form D round in the record; their debut lead averages $-0.22$ standard
deviations (SE $0.04$) against $-0.07$ (SE $0.04$) for the 910 funded,
29\% sit in the most lagging fifth against 22\%, and both groups are
mildly atypical ($+0.13$ and $+0.15$). Rounds before 2009, foreign
placements, and registered offerings post no Form D.

Value tiers: ten linked owners hold 18.5\% of \$19.2 trillion of linked
peak revenue, fifty hold 41.5\%, a hundred 53.6\%, five hundred 82.9\%.
In standard deviations of each axis across all debuts, with standard
errors: the ten largest sit at $-0.35$ ($0.19$) on atypicality and
$-0.01$ ($0.19$) on lead; the fifty largest at $-0.08$ ($0.13$) and
$+0.13$ ($0.16$), with 28\% in the most leading fifth; the hundred
largest at $-0.02$ ($0.09$) and $+0.02$ ($0.11$), with 25\% in the most
leading fifth; the five hundred largest at $+0.11$ ($0.05$) and $-0.04$
($0.05$), with 20\% in the most leading fifth. All 4{,}710 linked owners
average $+0.08$ on atypicality and $-0.05$ on lead.

Curated vocabularies: internet wording appears in 2.94 million
class-filing rows, AI in 167 thousand, blockchain in 114 thousand;
registration rates are 52.1\%, 34.5\% and 31.2\% against 57.6\% for all
debut filings (the last two cohorts are young). On the 2009--2018 ladder,
AI debuts are funded at 7.8\% and blockchain at 8.0\% against 2.8\% for
internet and 1.6\% overall, while at the IPO rung the gap disappears:
0.16\% of AI debuts, 0.22\% of blockchain, and 0.21\% of internet debuts
against 0.17\% overall (1{,}896 AI debuts, 1{,}366 blockchain).
